\documentclass{article}

\usepackage{polyglossia}
\setmainlanguage{lithuanian}
\usepackage{fontspec}
\setmainfont{TeX Gyre Termes}
\usepackage{indentfirst}
\frenchspacing

\usepackage[backend=biber]{biblatex}
\addbibresource{Bibliography.bib}

\usepackage{amsmath}
\usepackage{amsthm}
\usepackage{mathtools}
\usepackage{unicode-math}
\usepackage{physics}
\usepackage{upgreek}
\usepackage{siunitx} % dydžių vienetams
\sisetup{
output-decimal-marker = {,},
exponent-product = \cdot,
inter-unit-product = \ensuremath{\cdot} }

\theoremstyle{definition}
\newtheorem{definition}{Apibrėžimas}[section]

\theoremstyle{plain}
\newtheorem{preposition}{Teiginys}[section]

\theoremstyle{remark}
\newtheorem*{remark}{Pavyzdys}



\usepackage[
	pdfusetitle,
	bookmarks=true,
	bookmarksnumbered=true,
	bookmarksopen=false,
	breaklinks=false,
	backref=false,
	colorlinks=false
]{hyperref}
\hypersetup{
	pdfstartview=FitH,
	bookmarksopen=true,
	hidelinks,  
	pdfdisplaydoctitle=true,
	pdfpagemode=UseNone
}
\hypersetup{hidelinks}
\usepackage[a4paper]{geometry}
\geometry{tmargin=2cm,bmargin=2cm,lmargin=3cm,rmargin=1.5cm}
\usepackage{fancyhdr}
\pagestyle{fancy}

\usepackage{booktabs}
\usepackage{tocloft}
\usepackage{graphicx}

\renewcommand{\headrulewidth}{0pt}

\title{Taškinių grupių dvigubos grupės ir jų neredukuojami įvaizdžiai}

\begin{document}
\newcommand{\set}[1]{ \mathbb{#1} }
\newcommand{\spin}{ \mathup{Spin}(3) }
\newcommand{\pinp}{ \mathup{Pin}^{+}(3) }
\newcommand{\pinm}{ \mathup{Pin}^{-}(3) }
\newcommand{\pinpm}{ \mathup{Pin}^\pm(3) }

\global\long\def\imath{\mathrm{i}}%
 
\global\long\def\emath{\mathrm{e}}%
 
\rhead{}

\lhead{}

\cfoot{\thepage}

\normalfont

\thispagestyle{empty}
\begin{center}
Vilniaus universitetas\\
Fizikos fakultetas\\
Teorinės fizikos ir astrofizikos institutas
\par\end{center}

\vspace{3.5cm}

\begin{center}
	Martynas Lašas
\par\end{center}

\begin{center}
\textbf{TAŠKINIŲ GRUPIŲ DVIGUBOS GRUPĖS IR JŲ NEREDUKUOJAMI ĮVAIZDŽIAI}\vspace{0.8cm}
\par\end{center}

\begin{center}
BAKALAURO STUDIJŲ BAIGIAMASIS DARBAS 
\par\end{center}

\begin{center}
\vspace{0.8cm}
Fizikos studijų programa
\par\end{center}

\begin{center}
\vspace{3.5cm}
\par\end{center}

\begin{tabular*}{0.9\textwidth}{@{\extracolsep{\fill}}ll}
Studentas & Martynas Lašas
\tabularnewline[0.5cm]
	Leista ginti & 2023-05-20
	\tabularnewline
Darbo vadovas & dr. Vidas Regelskis
\tabularnewline[0.5cm]
	Instituto direktorius & Prof. Dr. Egidijus Anisimovas
	\tabularnewline
\end{tabular*}
\begin{center}
\vfill{}
\par\end{center}

\begin{center}
Vilnius 2025
\par\end{center}

\newpage{}

\normalfont\tableofcontents{}

\newpage{}

\section*{Įvadas\addcontentsline{toc}{section}{Įvadas}}

%
\section{Sukimų grupė ir jos apibendrinimai}

\subsection{Clifford'o algebros}


Toliau nagrinėsime nagrinėsime dvi Clifford'o algebras $\Clp$ ir $\Clm$.
Tai yra vektorinės erdvės su papildoma bitiesine daugybos operacija,
kuri leidžia sudauginti du vektorius ir gauti naują vektorių.
Šias algebras apibrėšime pirma paimdami daug bendresnę tenzorinę algebrą
ir tada jai suteiksime norimą struktūrą.

\begin{definition}
	Vienetinė tenzorinė algebra $T(\set R^3)$ yra tiesioginė suma
	\begin{equation}
		T(\mathbb{R}^3) 
		=
		\bigoplus_{k=0}^{\infty} T^k
		=
		\mathbb{R} 
		\oplus \set R^3 
		\oplus (\set R^3 \ot \set R^3) 
		\oplus (\set R^3 \ot \set R^3 \ot \set R^3) 
		\oplus \dots
		\,,
	\end{equation}
	$T^k$ čia yra vektorinės erdvės $T^0 = \set R$\: ,\: $T^1=\set R^3$\: ,\: $T^2=\set R^3 \ot \set R^3$\: ,\:
	$T^3=\set R^3 \ot \set R^3 \ot \set R^3$
	ir taip toliau.
	Algebros daugyba yra duodama vektorių tenzorinės sandaugos.
	Tegul $\mathbb1 \in T^0$ žymi vienetinį tenzorinės algebros elementą,
	kuris su visais $v \in T(\set R^3)$ tenkina
	$\mathbb1 \otimes v = v \otimes \mathbb1 = v$.
\end{definition}
Kitaip pasakius tai yra algebra,
kurios vektoriai yra įvairių ilgiu tenzorinės sandaugos
$\vb v_1 \ot \vb v_2 \ot \dots \ot \vb v_n$\:,\: $\vb v_i \in \set R^3$
ir tokių sandaugų tiesinės kombinacijos.
\begin{example}
	Vektorius $3\;\mathbb 1 + \vb v$ priklauso $T(\set R^3)$
	ir yra suma vektorių iš $T^0$ ir $T^1$.
	Vektorius $2\vb w + \vb v \ot \vb w$ yra suma vektorių iš $T^1$ ir $T^2$.
	Jų sandauga tenzorinėje algebroje yra
	\begin{equation}
		\begin{split}
			\qty(3\;\mathbb 1 + \vb v)
			\ot
			\qty(2\vb w + \vb v \ot \vb w)
			=&
			3\;\mathbb 1 \ot 2\vb w
			+ 3\;\mathbb 1 \ot \vb v \ot \vb w
			+ \vb v \ot 2\vb w
			+ \vb v \ot \vb v \ot \vb w
			\\=&
			3\cdot2\qty(\mathbb 1 \ot \vb w)
			+ 3\qty(\mathbb 1 \ot \vb v \ot \vb w)
			+ 2(\vb v \ot \vb w)
			+ \qty(\vb v \ot \vb v \ot \vb w)
			\\=&
			6 \vb w
			+ 3\qty(\vb v \ot \vb w)
			+ 2(\vb v \ot \vb w)
			+ \qty(\vb v \ot \vb v \ot \vb w)
			\\=&
			6 \vb w
			+ 5\qty(\vb v \ot \vb w)
			+ \qty(\vb v \ot \vb v \ot \vb w)
			\,,
		\end{split}
	\end{equation}
	čia pirmoje eilutėje pasinaudojome tuo kad tenzorinės sandauga
	distributyvi; antroje eilutėje pasinaudojame tuo kad tenzorinė sandauga
	tiesiška kiekvieno dauginamo vektoriaus atžvilgiu, kad iškeltume 
	skaliarinius daugiklius; trečioje eilutėje sudedame tiesiškai
	priklausomus vektorius.
	Gautas vektorius yra suma vektorių iš $T^1$, $T^2$ ir $T^3$.
\end{example}


Dabar $T(\set R^3)$ suteiksime papildomą struktūrą, kad algebroje
$\Clp$ būtų tenkinamas sąryšis
\begin{equation}
	\vb v \ot \vb w
	+
	\vb w \ot \vb v
	=
	2(\vb v , \vb w) \mathbb 1
	\quad ,
	\label{eq:cl+}
\end{equation}
o algebroje $\Clm$ tenkinamas sąryšis
\begin{equation}
	\vb v \ot \vb w
	+
	\vb w \ot \vb v
	=
	- 2(\vb v , \vb w) \mathbb 1
	\quad .
	\label{eq:cl-}
\end{equation}
Šiuos sąryšius ,,užkoduojame'' į aibes $S^\pm \subset T(\set R^3)$
\begin{equation}
	\begin{split}
		S^+ = \qty{
			\vb v \ot \vb w
			+
			\vb w \ot \vb v
			-
			2(\vb v , \vb w) \mathbb 1
			\mid
			\vb v, \vb w \in \set R^3
		}
		\\
		S^- = \qty{
			\vb v \ot \vb w
			+
			\vb w \ot \vb v
			+
			2(\vb v , \vb w) \mathbb 1
			\mid
			\vb v, \vb w \in \set R^3
		}\,.
	\end{split}
\end{equation}
Užrašome tenzorinės algebros poalgebrius $I^\pm \subset T(\set R^3)$,
gaunamus iš visų vektorių, kurie turi bent vieną daugiklį iš $S^\pm$
\begin{equation}
	I^\pm = \qty{
		\sum_{\substack{
			s \in S^\pm \\ 
			a,b \in \mathup{T(V)} 
		}}
		a \ot s \ot b
	}
	\quad .
\end{equation}
Šie poalgebriai papildomai tenkina savybę, kad bet kokiems $a \in T(\set R^3)$
ir $d \in I^\pm$, $a\ot d, d \ot a \in I^\pm$. Poalgebris
tenkinantis šia savybę yra vadinamas algebros idealu.
Literatūroje šią konstrukciją yra įprasta pateikti sakiniu:
,,$I^\pm$ yra algebros $T(\set R^3)$ idealas generuojamas 
sąryšio (\ref{eq:cl+}) ar atitinkamai (\ref{eq:cl-})''.

Duotam vektoriui $a \in T(\set R^3)$ aibė
$\qty{a+d \mid d \in I^\pm}$ (dar žymima $a+I^\pm$)
vadinama gretutine klase (ang. \emph{coset} arba \emph{equivalence class}),
o vektorius $a$ vadinamas klasės atstovu.
Dviejų gretutinių klasių $a+I^\pm$ ir $b + I^\pm$ sandauga yra gretutinė
klasė, kuriai priklauso atstovų sandauga $a \ot b$.
Ši sandauga nepriklauso nuo pasirinktų klasės atstovų, nes
\begin{equation}
	(a+I^\pm)\ot(b+I^\pm)
	=
	a \ot b
	+ a \ot I^\pm
	+ I^\pm \ot b
	+ I^\pm \ot I^\pm
	=
	a \ot b + I^\pm
	\,.
\end{equation}
Sumos nariai 
$a \ot I^\pm = I^\pm $,
$I^\pm \ot b = I^\pm $,
$I^\pm \ot I^\pm = I^\pm $,
nes $I^\pm$ yra algebros $T(\set R^3)$ idealas.
Tai tikrai
\begin{equation}
	(a+I^\pm)\ot(b+I^\pm)
	=
	a \ot b + I^\pm
	\,.
\end{equation}



\begin{definition}
	Algebra $\Clpm$ yra faktorinė algebra $T(\set R^3)/I^\pm$,
	kurios elementai yra gretutinės klasės $a+I^\pm$, $a \in T(\set R^3)$
	ir jų sandauga yra $(a+I^\pm)\ot(b+I^\pm)=a\ot b + I^\pm.$
\end{definition}

\begin{example}
	$\vb w \ot \vb v +\vb v \ot \vb w + I^+$ yra gretutinė klasė,
	kuri yra vektorius algebroje $\Clp$. 
	Supaprastiname išraišką
	\begin{equation}
		\vb w \ot \vb v +\vb v \ot \vb w + I^+
		= 
		\vb w \ot \vb v +\vb v \ot \vb w 
		- \qty(
			\vb v \ot \vb w + \vb w \ot \vb v 
			- 2(\vb v,\vb w)\mathbb 1
		) 
		+ I^+
		=
		2(\vb v,\vb w)\mathbb 1 + I^+
		\,,
	\end{equation}
	čia pasinaudojome tuo kad 
	$\qty( \vb v \ot \vb w + \vb w \ot \vb v - 2(\vb v,\vb w)\mathbb 1) \in I^+$.
	Pavyko parodyti, kad $\vb w \ot \vb v + \vb v \ot \vb w$ 
	ir $2(\vb v,\vb w) \mathbb 1$ atstovauja tą pačią klasę.
\end{example}

\begin{example}
	Tegul $a,b \in T(\set R^3)$ ir
	$ a \ot (\vb v \ot \vb w + 2(\vb v,\vb w)\mathbb 1) \ot b + I^- $
	yra gretutinė klasė, kuri yra vektorius iš $\Clm$. Šią išraišką supaprastiname atlikdami veiksmus
	\begin{equation}
		\begin{aligned}
			a \ot (\vb v \ot \vb w + 2(\vb v,\vb w)\mathbb 1) \ot b + I^- 
			& = a \ot (\vb v \ot \vb w + 2(\vb v,\vb w)\mathbb 1) \ot b
			\\& \quad - a \ot \qty(\vb v \ot \vb w + \vb w \ot \vb v +2(\vb v,\vb w) \mathbb 1)  \ot b + I^-
			\\& = a \ot \qty( \vb v \ot \vb w + 2(\vb v,\vb w)\mathbb 1 - \vb v \ot \vb w - \vb w \ot \vb v -2(\vb v,\vb w) \mathbb 1))\ot b + I^-
			\\& = a \ot \qty( - \vb w \ot \vb v )\ot b + I^-
			\\& = - a \ot	\vb w \ot \vb v	\ot b + I^- 
			\,.
		\end{aligned}
	\end{equation}
	Parodėme, kad sudarius faktorinę algebrą $\Clm$,
	galime jos viduje taikyti sąryšį (\ref{eq:cl-})
	manipuliuojant gretutinės klasės atstovus.
\end{example}
Toliau priimame sutrumpintą žymėjimą, nes išreikštai
su gretutinėmis klasėmis nebedirbsime. Užrašysime tik
klasės atstovą, o kad išraiška skirtųsi nuo 
vektorių tenzorinės sandaugos vietoje ženklo ,,$\ot$''
tiesiog rašysime vektorius vieną šalia kito.
Vietoje
\begin{equation}
	a \ot (\vb v \ot \vb w + \vb v \ot \vb w) \ot b + I^-
	=
	-2(\vb v,\vb w) \qty(a \ot b) + I^-
\end{equation}
rašysime
\begin{equation}
	a (\vb v \vb w + \vb v \vb w) b
	=
	-2(\vb v,\vb w) a b
	\,.
\end{equation}

\begin{proposition}
	\label{prop:R_action}
	Galioja tapatybės:
	\begin{enumerate}
		\item[$(i)$] algebrose $\Clpm$ turime $ \vb v \vb v = \pm(\vb v,\vb v) \mathbb 1$
		\item[$(ii)$] algebroje $\Clp$ turime $ -\vb v \vb w \vb v = -2(\vb v,\vb w) \vb v + (\vb v,\vb v)\vb w $
		\item[$(iii)$] algebroje $\Clm$ turime $ \vb v \vb w \vb v = -2(\vb v,\vb w) \vb v + (\vb v,\vb v)\vb w $
	\end{enumerate}
\end{proposition}
\begin{proof} Taikydami sąryšį $\vb v \vb w + \vb w \vb v = \pm2(\vb v ,\vb w)$ gauname
	\begin{enumerate}
		\item[$(i)$] 
			$
			\vb v \vb v 
			= \frac{1}{2}(\vb v \vb v + \vb v \vb v) 
			= \pm(\vb v,\vb v) \mathbb 1
			\,,
			$
		\item[$(ii)$] 	
			$
			-\vb v \vb w \vb v 
			= -\qty( 2(\vb v,\vb w)\mathbb 1 - \vb{wv} )\vb v 
			= -2(\vb v,\vb w) \vb v + \vb{wvv} 
			= -2(\vb v,\vb w) \vb v + (\vb v,\vb v)\vb w 
			\,,
			$
		\item[$(iii)$]
			$
			\vb v \vb w \vb v 
			= \qty( - 2(\vb v,\vb w)\mathbb 1 - \vb{wv} )\vb v 
			= -2(\vb v,\vb w) \vb v - \vb{wvv} 
			= -2(\vb v,\vb w) \vb v + (\vb v,\vb v)\vb w 
			\,.
			$
			\qedhere
	\end{enumerate}
\end{proof}

\begin{remark}
	Atkreipiame dėmesį, kad kai $(\vb v,\vb v)=1$,
	$(\vb v,\vb w)\vb v$ yra vektoriaus $\vb w$ projekcija į $\vb v$
	ir $-2(\vb v,\vb w)\vb v + \vb w$ yra vektoriaus atspindys
	per plokštumą statmeną vektoriui $\vb v$. 
\end{remark}

Tegul $\qty{\vb e_1,\vb e_2,\vb e_3}$ yra ortonormali $\set R^3$ bazė.
Bazinius vektorius galime perkelti į algebra $\mathup{Cl^\pm(V)}$,
kur jie turi tenkinti
\begin{equation}
	\begin{split}
		\vb e_i \vb e_j &= - \vb e_j \vb e_i \ , \ i \neq j \\
		\vb e_i \vb e_i &= \pm \mathbb 1
		\quad .
	\end{split}
	\label{eq:base-relations}
\end{equation}
Su šių sąryšių pagalba galime bet kokį $\mathup{Cl^\pm(V)}$ elementą užrašyti,
kaip tiesinę kombinaciją aštuonių elementų:
\begin{equation}
	\mathbb 1 \,,\quad
	\vb e_1 \,,\quad 
	\vb e_2 \,,\quad  
	\vb e_3 \,,\quad  
	\vb e_2 \vb e_3 \,,\quad
	\vb e_1 \vb e_3 \,,\quad
	\vb e_1 \vb e_2 \,,\quad 
	\vb e_1 \vb e_2 \vb e_3 \,,
\end{equation}
nes su pirma taisykle bet kurį tenzorinės algebros bazės elementą galima išrikiuoti,
taip kad ortų indeksai būtų surašyti didėjimo tvarka
ir su antra taisykle pasikartojantys ortai panaikinami.
Tai reiškia $\Clpm$ yra aštuonmetės algebros.
Įvedame pažymėjimą, kurį toliau naudosime algebrų $\Clpm$ bazei
\begin{equation}
	\vb e_{23} := \vb e_2 \vb e_3
	\,,\quad
	\vb e_{31} := \vb e_3 \vb e_1 = -\vb e_1 \vb e_3 
	\,,\quad
	\vb e_{12} := \vb e_1 \vb e_2
	\,,\quad
	\vb e_{123} := \vb e_1 \vb e_2 \vb e_3
	\,.
\end{equation}

\begin{definition}
	Transponuota $\Clpm$ bazė yra 
	\begin{equation}
		\qty(\vb e_{i_1} \vb e_{i_2} \dots \vb e_{i_k} )^t
		:=
		\vb e_{i_k} \vb e_{i_{k-1}} \dots \vb e_{i_1}
		=
		(-1)^\frac{k(k-1)}{2} \vb e_{i_1} \vb e_{i_2} \dots \vb e_{i_k} 
	\end{equation}
	ir transponavimo operacija yra tiesiškai praplečiama
	bet kuriam algebros $\Clpm$ vektoriui, veikimu ant bazės.
	Transponavimas yra algebros anti-automorfizmas,
	tai reiškia, kad bet kuriems algebros elementams $a$ ir $b$ galioja
	\begin{equation}
		(a b)^t = b^t a^t
		\,.
	\end{equation}
\end{definition}
\begin{definition}
	Apibrėžiame algebrų $\Clpm$ automorfizmą $\alpha$, jo veikimu ant algebrų bazės
	\begin{equation}
		\alpha( \vb e_{i_1} \vb e_{i_2} \dots \vb e_{i_k} )
		:=
		(-1)^{k} \vb e_{i_1} \vb e_{i_2} \dots \vb e_{i_k} 
		\,.
	\end{equation}
\end{definition}
Nesunku pamatyti,
kad nėra svarbu kokia tvarka pritaikomas transponavimas ir automorfizmas $\alpha$,
tai $\alpha(v)^t = \alpha(v^t)$.


Pažymime algebros poerdvius
\begin{equation}
	\begin{split}
		A_0^\pm &:= \qty{
			v \in \Clpm \mid \alpha(v)=v
		}
		\,,
		\\
		A_1^\pm &:= \qty{
			v \in \Clpm \mid \alpha(v)=-v
		}
		\,.
	\end{split}
\end{equation}
Susiejant su erdvės baze šiuos poerdvius galima užrašyti kaip
tiesinį apvalką pasirinktų bazės elementų
\begin{equation}
	\begin{aligned}
		A_0^\pm &= \mathup{span}\qty(
			\mathbb 1,\,
			\vb e_{23},\,
			\vb e_{31},\,
			\vb e_{12}
		)
		\,,
		\\
		A_1^\pm &= \mathup{span}\qty(
			\vb e_1 ,\,
			\vb e_2 ,\,
			\vb e_3 ,\,
			\vb e_{123}
		)
		\,.
	\end{aligned}
\end{equation}

\begin{example}
	Tegul $v_0, v'_0 \in A_0^\pm$ ir $v_1, v'_1 \in A_1^\pm$. Turime
	\begin{equation}
		\begin{split}
			\alpha(v_0 v'_0)
			= \alpha(v_0)\alpha(v'_0)
			= v_0 v'_0
			\implies &
			v_0 v'_0 \in A_0^\pm
			\,,
			\\
			\alpha(v_0 v_1)
			= \alpha(v_0)\alpha(v_1)
			= v_0 (-v_1)
			= - v_0 v_1
			\implies &
			v_0 v_1 \in A_1^\pm
			\,,
			\\
			\alpha(v_1 v_0)
			= \alpha(v_1)\alpha(v_0)
			= (- v_1)v_0
			= - v_1 v_0
			\implies &
			v_1 v_0 \in A_1^\pm
			\,,
			\\
			\alpha(v_1 v'_1)
			= \alpha(v_1)\alpha(v'_1)
			= (- v_1)(- v'_1)
			= v_1 v'_1
			\implies &
			v_1 v'_1 \in A_0^\pm
			\,.
		\end{split}
	\end{equation}
	Taip yra tokia pati struktūra kaip grupė 
	$\set Z_2 \cong C_2$ su sudėtimi
	\begin{equation}
		\begin{split}
			\mathbf{lyginis} + \mathbf{lyginis} &= \mathbf{lyginis}
			\,,
			\\
			\mathbf{lyginis} + \mathbf{nelyginis} &= \mathbf{nelyginis}
			\,,
			\\
			\mathbf{nelyginis} + \mathbf{lyginis} &= \mathbf{nelyginis}
			\,,
			\\
			\mathbf{nelyginis} + \mathbf{nelyginis} &= \mathbf{lyginis}
			\,.
		\end{split}
	\end{equation}
\end{example}
Dėl tokios daugybos struktūros vektoriai erdvėje $A_0$ yra vadinami \emph{lyginiais},
o vektoriai erdvėje $A_1$ yra vadinami \emph{nelyginiais}.
Pamatome, kad automorfizmas $\alpha$ 
leidžia algebrą $\Clpm$ suskaidyti 
į lyginio ir nelyginio poerdvių tiesioginę sumą $A_0 \oplus A_1$.
Šis algebros išskaidymas į vektorinius poerdvius vadinamas
,,$\set Z_2$ gradavimu'', kurį suteikia automorfizmas $\alpha$.
Lyginis poerdvis yra uždaras algebros daugybos atžvilgiu,
todėl tai yra $\Clpm$ poalgebris.

\begin{proposition}
	\label{prop:even_iso}
	Algebrų $\Clp$ ir $\Clm$
	Lyginiai poalgebriai $A^+_0$ ir $A^-_0$ yra izomorfiškos algebros.
\end{proposition}

\begin{proof}
Algebros $A_0^+$ ir $A_0^-$ turi bazę $\qty{
		\mathbb 1,\, 
		\vb e_{23},\,
		\vb e_{31},\,
		\vb e_{12}
	}$. 
Pastebime kad algebroje $A_0^+$
\begin{gather}
	\vb e_{23} \vb e_{31} 
	= \vb e_2 \vb e_3 \vb e_3 \vb e_1
	= \vb e_2 \vb e_1
	= - \vb e_{12}
	\,,
\\
	\vb e_{31} \vb e_{23} 
	= \vb e_3 \vb e_1 \vb e_2 \vb e_3
	= \vb e_1 \vb e_3 \vb e_3 \vb e_2
	= \vb e_1 \vb e_2
	= \vb e_{12}
	\,,
\\
	{\vb e_{31}}^2 
	= \vb e_3 \vb e_1 \vb e_3 \vb e_1 
	= - \vb e_3 \vb e_3 \vb e_1 \vb e_1 
	= - \mathbb 1
	\,,
\\
	{\vb e_{23}}^2 
	= \vb e_2 \vb e_3 \vb e_2 \vb e_3 
	= - \vb e_2 \vb e_2 \vb e_3 \vb e_3 
	= - \mathbb 1
	\,,
\end{gather}
o algebroje $A_0^-$
\begin{gather}
	\vb e_{23} \vb e_{31} 
	= \vb e_2 \vb e_3 \vb e_3 \vb e_1 
	= - \vb e_2 \vb e_1 
	= \vb e_1 \vb e_2 
	= \vb e_{12}
	\,,
\\
	\vb e_{31} \vb e_{23} 
	= \vb e_3 \vb e_1 \vb e_2 \vb e_3
	= \vb e_1 \vb e_3 \vb e_3 \vb e_2
	= - \vb e_1 \vb e_2
	= - \vb e_{12}
	\,,
\\
	{\vb e_{31}}^2 
	= \vb e_3 \vb e_1 \vb e_3 \vb e_1 
	= - \vb e_3 \vb e_3 \vb e_1 \vb e_1 
	= - \mathbb 1
	\,,
\\
	{\vb e_{23}}^2 
	= \vb e_2 \vb e_3 \vb e_2 \vb e_3 
	= - \vb e_2 \vb e_2 \vb e_3 \vb e_3 
	= - \mathbb 1
	\,.
\end{gather}
Abiems algebroms $\vb e_{31}$ ir $\vb e_{23}$ yra algebros generatoriai.
Jie tenkina sąryšius
\begin{equation}
	\begin{aligned}
		\vb e_{31} \vb e_{23} &= - \vb e_{23} \vb e_{31} \,,\\
		{\vb e_{31}}^2 &= - \mathbb1 \,,\\
		{\vb e_{23}}^2 &= - \mathbb1 \,.
	\end{aligned}
\end{equation}

Tegul $f:A^+_0 \to A^-_0$ yra tiesinis atvaizdavimas,
apibrėžiamas veikimu ant $A^+_0$ bazės
\begin{equation}
	\mathbb1 \mapsto \mathbb1
	\,,\quad
	\vb e_{23} \mapsto -\vb e_{23}
	\,,\quad
	\vb e_{31} \mapsto -\vb e_{31}
	\,,\quad
	\vb e_{12} \mapsto -\vb e_{12}
	\,.
\end{equation}
Akivaizdu kad $f$ yra bijekcija.
Įrodymui kad $f$ yra algebrų homomorfizmas užtenka parodyti,
kad $f$ nepakeičia generatorių daugybos sąryšių
\begin{equation}
	\begin{aligned}
		f(\vb e_{31}) f(\vb e_{23}) 
		=(-\vb e_{31})(-\vb e_{23}) 
		&= - \vb e_{23} \vb e_{31} 
		= - f(\vb e_{23}) f(\vb e_{31}) 
\,,\\
		f(\vb e_{31})^2 
		=(-\vb e_{31})^2 
		&= - \mathbb1
		= - f(\mathbb1)
\,,\\
		f(\vb e_{23})^2 
		=(-\vb e_{23})^2 
		&= - \mathbb1
		= - f(\mathbb1)
\,.
	\end{aligned}
\end{equation}
Atvaizdavimas $f$ yra algebrų izomorfizmas,
tad $A^+_0 \cong A^-_0$.
\end{proof}


\subsection{Grupės $\pinp$ ir $\pinm$}

Algebros $\mathup{Cl^\pm(V)}$ poaibis sudarytas iš elemtų $\vb v_1 \vb v_2
\dots \vb v_N$, kur visi $\vb v_i$ tenkina $(\vb v_i,\vb v_i) = 1$, kartu su
algebros daugyba sudaro atitinkamą grupę $\pinpm$. Grupės uždarumo sąvybė akivaizdžiai tenkinama, 
vienetinis elementas yra tas pats kaip ir algebros vienetinis elementas $\mathbb 1$,
o kaip gaunamas atvikštinis elementas siriasi grupeių $\pinp$ ir $\pinm$.

Grupėje $\pinp$ elemento $v$ atvirkštinis elementas yra tranponuotas $v^t$
\begin{multline}
	\vb v_1 \vb v_2 \dots \vb v_N (\vb v_1 \vb v_2 \dots \vb v_N )^t
	=
	\vb v_1 \vb v_2 \dots \vb v_N \vb v_N \vb v_{N-1} \dots \vb v_1
	=
	(\vb v_N,\vb v_N)
	(\vb v_{N-1},\vb v_{N-1})
	\dots
	(\vb v_1,\vb v_1)
	\mathbb 1
	=
	\mathbb 1
\end{multline}
Grupėje $\pinm$ atvirkštinis elementas yra $\alpha(v^t)$
\begin{multline}
	\vb v_1 \vb v_2 \dots \vb v_N \alpha\qty(\vb v_1 \vb v_2 \dots \vb v_N )^t
	=
	\vb v_1 \vb v_2 \dots \vb v_N \alpha(\vb v_N) \alpha(\vb v_{N-1}) \dots \alpha(\vb v_1)
	\\=
	\vb v_1 \vb v_2 \dots \vb v_N (-\vb v_N) (-\vb v_{N-1}) \dots (-\vb v_1)
	=
	(\vb v_N,\vb v_N)
	(\vb v_{N-1},\vb v_{N-1})
	\dots
	(\vb v_1,\vb v_1)
	\mathbb 1
	=
	\mathbb 1
\end{multline}

Grupė $\spin$ galima gauti dviem būdais: kaip $\pinp$ pogrupį sudarytą
iš lyginių vektorių sandaugų arba kaip $\pinm$ pogrupį gautą taip pat.
Abi konstrukcijos yra viena kitai izomofiškos, nes minusas iš sandaugos
$\pinm$ pasirodys lyginį kartą sakičių ir rezultatas bus toks pat kaip
$\pinp$ pogrupyje.

\subsection{$\spin$ ir $\pinpm$ veikimas ant $\set R^3$}

Visos aptartos grupės veikia vektorius iš $\set R^3$, juos pirma patalpinant į
$\mathup{Cl^\pm(V)}$ ir tada su algebros sandauga
\begin{equation}
	\rho(v).x = \begin{cases}
		vxv^t		& v \in \pinp \\
		\alpha(v)xv^t	& v \in \pinm 
	\end{cases}
\end{equation}






%\clearpage{}


\section{Sukimų grupė ir jos apibendrinimai}


\section{Sukimų grupė ir jos apibendrinimai}

\subsection{Clifford'o algebros}


Toliau nagrinėsime nagrinėsime dvi Clifford'o algebras $\Clp$ ir $\Clm$.
Tai yra vektorinės erdvės su papildoma bitiesine daugybos operacija,
kuri leidžia sudauginti du vektorius ir gauti naują vektorių.
Šias algebras apibrėšime pirma paimdami daug bendresnę tenzorinę algebrą
ir tada jai suteiksime norimą struktūrą.

\begin{definition}
	Vienetinė tenzorinė algebra $T(\set R^3)$ yra tiesioginė suma
	\begin{equation}
		T(\mathbb{R}^3) 
		=
		\bigoplus_{k=0}^{\infty} T^k
		=
		\mathbb{R} 
		\oplus \set R^3 
		\oplus (\set R^3 \ot \set R^3) 
		\oplus (\set R^3 \ot \set R^3 \ot \set R^3) 
		\oplus \dots
		\,,
	\end{equation}
	$T^k$ čia yra vektorinės erdvės $T^0 = \set R$\: ,\: $T^1=\set R^3$\: ,\: $T^2=\set R^3 \ot \set R^3$\: ,\:
	$T^3=\set R^3 \ot \set R^3 \ot \set R^3$
	ir taip toliau.
	Algebros daugyba yra duodama vektorių tenzorinės sandaugos.
	Tegul $\mathbb1 \in T^0$ žymi vienetinį tenzorinės algebros elementą,
	kuris su visais $v \in T(\set R^3)$ tenkina
	$\mathbb1 \otimes v = v \otimes \mathbb1 = v$.
\end{definition}
Kitaip pasakius tai yra algebra,
kurios vektoriai yra įvairių ilgiu tenzorinės sandaugos
$\vb v_1 \ot \vb v_2 \ot \dots \ot \vb v_n$\:,\: $\vb v_i \in \set R^3$
ir tokių sandaugų tiesinės kombinacijos.
\begin{example}
	Vektorius $3\;\mathbb 1 + \vb v$ priklauso $T(\set R^3)$
	ir yra suma vektorių iš $T^0$ ir $T^1$.
	Vektorius $2\vb w + \vb v \ot \vb w$ yra suma vektorių iš $T^1$ ir $T^2$.
	Jų sandauga tenzorinėje algebroje yra
	\begin{equation}
		\begin{split}
			\qty(3\;\mathbb 1 + \vb v)
			\ot
			\qty(2\vb w + \vb v \ot \vb w)
			=&
			3\;\mathbb 1 \ot 2\vb w
			+ 3\;\mathbb 1 \ot \vb v \ot \vb w
			+ \vb v \ot 2\vb w
			+ \vb v \ot \vb v \ot \vb w
			\\=&
			3\cdot2\qty(\mathbb 1 \ot \vb w)
			+ 3\qty(\mathbb 1 \ot \vb v \ot \vb w)
			+ 2(\vb v \ot \vb w)
			+ \qty(\vb v \ot \vb v \ot \vb w)
			\\=&
			6 \vb w
			+ 3\qty(\vb v \ot \vb w)
			+ 2(\vb v \ot \vb w)
			+ \qty(\vb v \ot \vb v \ot \vb w)
			\\=&
			6 \vb w
			+ 5\qty(\vb v \ot \vb w)
			+ \qty(\vb v \ot \vb v \ot \vb w)
			\,,
		\end{split}
	\end{equation}
	čia pirmoje eilutėje pasinaudojome tuo kad tenzorinės sandauga
	distributyvi; antroje eilutėje pasinaudojame tuo kad tenzorinė sandauga
	tiesiška kiekvieno dauginamo vektoriaus atžvilgiu, kad iškeltume 
	skaliarinius daugiklius; trečioje eilutėje sudedame tiesiškai
	priklausomus vektorius.
	Gautas vektorius yra suma vektorių iš $T^1$, $T^2$ ir $T^3$.
\end{example}


Dabar $T(\set R^3)$ suteiksime papildomą struktūrą, kad algebroje
$\Clp$ būtų tenkinamas sąryšis
\begin{equation}
	\vb v \ot \vb w
	+
	\vb w \ot \vb v
	=
	2(\vb v , \vb w) \mathbb 1
	\quad ,
	\label{eq:cl+}
\end{equation}
o algebroje $\Clm$ tenkinamas sąryšis
\begin{equation}
	\vb v \ot \vb w
	+
	\vb w \ot \vb v
	=
	- 2(\vb v , \vb w) \mathbb 1
	\quad .
	\label{eq:cl-}
\end{equation}
Šiuos sąryšius ,,užkoduojame'' į aibes $S^\pm \subset T(\set R^3)$
\begin{equation}
	\begin{split}
		S^+ = \qty{
			\vb v \ot \vb w
			+
			\vb w \ot \vb v
			-
			2(\vb v , \vb w) \mathbb 1
			\mid
			\vb v, \vb w \in \set R^3
		}
		\\
		S^- = \qty{
			\vb v \ot \vb w
			+
			\vb w \ot \vb v
			+
			2(\vb v , \vb w) \mathbb 1
			\mid
			\vb v, \vb w \in \set R^3
		}\,.
	\end{split}
\end{equation}
Užrašome tenzorinės algebros poalgebrius $I^\pm \subset T(\set R^3)$,
gaunamus iš visų vektorių, kurie turi bent vieną daugiklį iš $S^\pm$
\begin{equation}
	I^\pm = \qty{
		\sum_{\substack{
			s \in S^\pm \\ 
			a,b \in \mathup{T(V)} 
		}}
		a \ot s \ot b
	}
	\quad .
\end{equation}
Šie poalgebriai papildomai tenkina savybę, kad bet kokiems $a \in T(\set R^3)$
ir $d \in I^\pm$, $a\ot d, d \ot a \in I^\pm$. Poalgebris
tenkinantis šia savybę yra vadinamas algebros idealu.
Literatūroje šią konstrukciją yra įprasta pateikti sakiniu:
,,$I^\pm$ yra algebros $T(\set R^3)$ idealas generuojamas 
sąryšio (\ref{eq:cl+}) ar atitinkamai (\ref{eq:cl-})''.

Duotam vektoriui $a \in T(\set R^3)$ aibė
$\qty{a+d \mid d \in I^\pm}$ (dar žymima $a+I^\pm$)
vadinama gretutine klase (ang. \emph{coset} arba \emph{equivalence class}),
o vektorius $a$ vadinamas klasės atstovu.
Dviejų gretutinių klasių $a+I^\pm$ ir $b + I^\pm$ sandauga yra gretutinė
klasė, kuriai priklauso atstovų sandauga $a \ot b$.
Ši sandauga nepriklauso nuo pasirinktų klasės atstovų, nes
\begin{equation}
	(a+I^\pm)\ot(b+I^\pm)
	=
	a \ot b
	+ a \ot I^\pm
	+ I^\pm \ot b
	+ I^\pm \ot I^\pm
	=
	a \ot b + I^\pm
	\,.
\end{equation}
Sumos nariai 
$a \ot I^\pm = I^\pm $,
$I^\pm \ot b = I^\pm $,
$I^\pm \ot I^\pm = I^\pm $,
nes $I^\pm$ yra algebros $T(\set R^3)$ idealas.
Tai tikrai
\begin{equation}
	(a+I^\pm)\ot(b+I^\pm)
	=
	a \ot b + I^\pm
	\,.
\end{equation}



\begin{definition}
	Algebra $\Clpm$ yra faktorinė algebra $T(\set R^3)/I^\pm$,
	kurios elementai yra gretutinės klasės $a+I^\pm$, $a \in T(\set R^3)$
	ir jų sandauga yra $(a+I^\pm)\ot(b+I^\pm)=a\ot b + I^\pm.$
\end{definition}

\begin{example}
	$\vb w \ot \vb v +\vb v \ot \vb w + I^+$ yra gretutinė klasė,
	kuri yra vektorius algebroje $\Clp$. 
	Supaprastiname išraišką
	\begin{equation}
		\vb w \ot \vb v +\vb v \ot \vb w + I^+
		= 
		\vb w \ot \vb v +\vb v \ot \vb w 
		- \qty(
			\vb v \ot \vb w + \vb w \ot \vb v 
			- 2(\vb v,\vb w)\mathbb 1
		) 
		+ I^+
		=
		2(\vb v,\vb w)\mathbb 1 + I^+
		\,,
	\end{equation}
	čia pasinaudojome tuo kad 
	$\qty( \vb v \ot \vb w + \vb w \ot \vb v - 2(\vb v,\vb w)\mathbb 1) \in I^+$.
	Pavyko parodyti, kad $\vb w \ot \vb v + \vb v \ot \vb w$ 
	ir $2(\vb v,\vb w) \mathbb 1$ atstovauja tą pačią klasę.
\end{example}

\begin{example}
	Tegul $a,b \in T(\set R^3)$ ir
	$ a \ot (\vb v \ot \vb w + 2(\vb v,\vb w)\mathbb 1) \ot b + I^- $
	yra gretutinė klasė, kuri yra vektorius iš $\Clm$. Šią išraišką supaprastiname atlikdami veiksmus
	\begin{equation}
		\begin{aligned}
			a \ot (\vb v \ot \vb w + 2(\vb v,\vb w)\mathbb 1) \ot b + I^- 
			& = a \ot (\vb v \ot \vb w + 2(\vb v,\vb w)\mathbb 1) \ot b
			\\& \quad - a \ot \qty(\vb v \ot \vb w + \vb w \ot \vb v +2(\vb v,\vb w) \mathbb 1)  \ot b + I^-
			\\& = a \ot \qty( \vb v \ot \vb w + 2(\vb v,\vb w)\mathbb 1 - \vb v \ot \vb w - \vb w \ot \vb v -2(\vb v,\vb w) \mathbb 1))\ot b + I^-
			\\& = a \ot \qty( - \vb w \ot \vb v )\ot b + I^-
			\\& = - a \ot	\vb w \ot \vb v	\ot b + I^- 
			\,.
		\end{aligned}
	\end{equation}
	Parodėme, kad sudarius faktorinę algebrą $\Clm$,
	galime jos viduje taikyti sąryšį (\ref{eq:cl-})
	manipuliuojant gretutinės klasės atstovus.
\end{example}
Toliau priimame sutrumpintą žymėjimą, nes išreikštai
su gretutinėmis klasėmis nebedirbsime. Užrašysime tik
klasės atstovą, o kad išraiška skirtųsi nuo 
vektorių tenzorinės sandaugos vietoje ženklo ,,$\ot$''
tiesiog rašysime vektorius vieną šalia kito.
Vietoje
\begin{equation}
	a \ot (\vb v \ot \vb w + \vb v \ot \vb w) \ot b + I^-
	=
	-2(\vb v,\vb w) \qty(a \ot b) + I^-
\end{equation}
rašysime
\begin{equation}
	a (\vb v \vb w + \vb v \vb w) b
	=
	-2(\vb v,\vb w) a b
	\,.
\end{equation}

\begin{proposition}
	\label{prop:R_action}
	Galioja tapatybės:
	\begin{enumerate}
		\item[$(i)$] algebrose $\Clpm$ turime $ \vb v \vb v = \pm(\vb v,\vb v) \mathbb 1$
		\item[$(ii)$] algebroje $\Clp$ turime $ -\vb v \vb w \vb v = -2(\vb v,\vb w) \vb v + (\vb v,\vb v)\vb w $
		\item[$(iii)$] algebroje $\Clm$ turime $ \vb v \vb w \vb v = -2(\vb v,\vb w) \vb v + (\vb v,\vb v)\vb w $
	\end{enumerate}
\end{proposition}
\begin{proof} Taikydami sąryšį $\vb v \vb w + \vb w \vb v = \pm2(\vb v ,\vb w)$ gauname
	\begin{enumerate}
		\item[$(i)$] 
			$
			\vb v \vb v 
			= \frac{1}{2}(\vb v \vb v + \vb v \vb v) 
			= \pm(\vb v,\vb v) \mathbb 1
			\,,
			$
		\item[$(ii)$] 	
			$
			-\vb v \vb w \vb v 
			= -\qty( 2(\vb v,\vb w)\mathbb 1 - \vb{wv} )\vb v 
			= -2(\vb v,\vb w) \vb v + \vb{wvv} 
			= -2(\vb v,\vb w) \vb v + (\vb v,\vb v)\vb w 
			\,,
			$
		\item[$(iii)$]
			$
			\vb v \vb w \vb v 
			= \qty( - 2(\vb v,\vb w)\mathbb 1 - \vb{wv} )\vb v 
			= -2(\vb v,\vb w) \vb v - \vb{wvv} 
			= -2(\vb v,\vb w) \vb v + (\vb v,\vb v)\vb w 
			\,.
			$
			\qedhere
	\end{enumerate}
\end{proof}

\begin{remark}
	Atkreipiame dėmesį, kad kai $(\vb v,\vb v)=1$,
	$(\vb v,\vb w)\vb v$ yra vektoriaus $\vb w$ projekcija į $\vb v$
	ir $-2(\vb v,\vb w)\vb v + \vb w$ yra vektoriaus atspindys
	per plokštumą statmeną vektoriui $\vb v$. 
\end{remark}

Tegul $\qty{\vb e_1,\vb e_2,\vb e_3}$ yra ortonormali $\set R^3$ bazė.
Bazinius vektorius galime perkelti į algebra $\mathup{Cl^\pm(V)}$,
kur jie turi tenkinti
\begin{equation}
	\begin{split}
		\vb e_i \vb e_j &= - \vb e_j \vb e_i \ , \ i \neq j \\
		\vb e_i \vb e_i &= \pm \mathbb 1
		\quad .
	\end{split}
	\label{eq:base-relations}
\end{equation}
Su šių sąryšių pagalba galime bet kokį $\mathup{Cl^\pm(V)}$ elementą užrašyti,
kaip tiesinę kombinaciją aštuonių elementų:
\begin{equation}
	\mathbb 1 \,,\quad
	\vb e_1 \,,\quad 
	\vb e_2 \,,\quad  
	\vb e_3 \,,\quad  
	\vb e_2 \vb e_3 \,,\quad
	\vb e_1 \vb e_3 \,,\quad
	\vb e_1 \vb e_2 \,,\quad 
	\vb e_1 \vb e_2 \vb e_3 \,,
\end{equation}
nes su pirma taisykle bet kurį tenzorinės algebros bazės elementą galima išrikiuoti,
taip kad ortų indeksai būtų surašyti didėjimo tvarka
ir su antra taisykle pasikartojantys ortai panaikinami.
Tai reiškia $\Clpm$ yra aštuonmetės algebros.
Įvedame pažymėjimą, kurį toliau naudosime algebrų $\Clpm$ bazei
\begin{equation}
	\vb e_{23} := \vb e_2 \vb e_3
	\,,\quad
	\vb e_{31} := \vb e_3 \vb e_1 = -\vb e_1 \vb e_3 
	\,,\quad
	\vb e_{12} := \vb e_1 \vb e_2
	\,,\quad
	\vb e_{123} := \vb e_1 \vb e_2 \vb e_3
	\,.
\end{equation}

\begin{definition}
	Transponuota $\Clpm$ bazė yra 
	\begin{equation}
		\qty(\vb e_{i_1} \vb e_{i_2} \dots \vb e_{i_k} )^t
		:=
		\vb e_{i_k} \vb e_{i_{k-1}} \dots \vb e_{i_1}
		=
		(-1)^\frac{k(k-1)}{2} \vb e_{i_1} \vb e_{i_2} \dots \vb e_{i_k} 
	\end{equation}
	ir transponavimo operacija yra tiesiškai praplečiama
	bet kuriam algebros $\Clpm$ vektoriui, veikimu ant bazės.
	Transponavimas yra algebros anti-automorfizmas,
	tai reiškia, kad bet kuriems algebros elementams $a$ ir $b$ galioja
	\begin{equation}
		(a b)^t = b^t a^t
		\,.
	\end{equation}
\end{definition}
\begin{definition}
	Apibrėžiame algebrų $\Clpm$ automorfizmą $\alpha$, jo veikimu ant algebrų bazės
	\begin{equation}
		\alpha( \vb e_{i_1} \vb e_{i_2} \dots \vb e_{i_k} )
		:=
		(-1)^{k} \vb e_{i_1} \vb e_{i_2} \dots \vb e_{i_k} 
		\,.
	\end{equation}
\end{definition}
Nesunku pamatyti,
kad nėra svarbu kokia tvarka pritaikomas transponavimas ir automorfizmas $\alpha$,
tai $\alpha(v)^t = \alpha(v^t)$.


Pažymime algebros poerdvius
\begin{equation}
	\begin{split}
		A_0^\pm &:= \qty{
			v \in \Clpm \mid \alpha(v)=v
		}
		\,,
		\\
		A_1^\pm &:= \qty{
			v \in \Clpm \mid \alpha(v)=-v
		}
		\,.
	\end{split}
\end{equation}
Susiejant su erdvės baze šiuos poerdvius galima užrašyti kaip
tiesinį apvalką pasirinktų bazės elementų
\begin{equation}
	\begin{aligned}
		A_0^\pm &= \mathup{span}\qty(
			\mathbb 1,\,
			\vb e_{23},\,
			\vb e_{31},\,
			\vb e_{12}
		)
		\,,
		\\
		A_1^\pm &= \mathup{span}\qty(
			\vb e_1 ,\,
			\vb e_2 ,\,
			\vb e_3 ,\,
			\vb e_{123}
		)
		\,.
	\end{aligned}
\end{equation}

\begin{example}
	Tegul $v_0, v'_0 \in A_0^\pm$ ir $v_1, v'_1 \in A_1^\pm$. Turime
	\begin{equation}
		\begin{split}
			\alpha(v_0 v'_0)
			= \alpha(v_0)\alpha(v'_0)
			= v_0 v'_0
			\implies &
			v_0 v'_0 \in A_0^\pm
			\,,
			\\
			\alpha(v_0 v_1)
			= \alpha(v_0)\alpha(v_1)
			= v_0 (-v_1)
			= - v_0 v_1
			\implies &
			v_0 v_1 \in A_1^\pm
			\,,
			\\
			\alpha(v_1 v_0)
			= \alpha(v_1)\alpha(v_0)
			= (- v_1)v_0
			= - v_1 v_0
			\implies &
			v_1 v_0 \in A_1^\pm
			\,,
			\\
			\alpha(v_1 v'_1)
			= \alpha(v_1)\alpha(v'_1)
			= (- v_1)(- v'_1)
			= v_1 v'_1
			\implies &
			v_1 v'_1 \in A_0^\pm
			\,.
		\end{split}
	\end{equation}
	Taip yra tokia pati struktūra kaip grupė 
	$\set Z_2 \cong C_2$ su sudėtimi
	\begin{equation}
		\begin{split}
			\mathbf{lyginis} + \mathbf{lyginis} &= \mathbf{lyginis}
			\,,
			\\
			\mathbf{lyginis} + \mathbf{nelyginis} &= \mathbf{nelyginis}
			\,,
			\\
			\mathbf{nelyginis} + \mathbf{lyginis} &= \mathbf{nelyginis}
			\,,
			\\
			\mathbf{nelyginis} + \mathbf{nelyginis} &= \mathbf{lyginis}
			\,.
		\end{split}
	\end{equation}
\end{example}
Dėl tokios daugybos struktūros vektoriai erdvėje $A_0$ yra vadinami \emph{lyginiais},
o vektoriai erdvėje $A_1$ yra vadinami \emph{nelyginiais}.
Pamatome, kad automorfizmas $\alpha$ 
leidžia algebrą $\Clpm$ suskaidyti 
į lyginio ir nelyginio poerdvių tiesioginę sumą $A_0 \oplus A_1$.
Šis algebros išskaidymas į vektorinius poerdvius vadinamas
,,$\set Z_2$ gradavimu'', kurį suteikia automorfizmas $\alpha$.
Lyginis poerdvis yra uždaras algebros daugybos atžvilgiu,
todėl tai yra $\Clpm$ poalgebris.

\begin{proposition}
	\label{prop:even_iso}
	Algebrų $\Clp$ ir $\Clm$
	Lyginiai poalgebriai $A^+_0$ ir $A^-_0$ yra izomorfiškos algebros.
\end{proposition}

\begin{proof}
Algebros $A_0^+$ ir $A_0^-$ turi bazę $\qty{
		\mathbb 1,\, 
		\vb e_{23},\,
		\vb e_{31},\,
		\vb e_{12}
	}$. 
Pastebime kad algebroje $A_0^+$
\begin{gather}
	\vb e_{23} \vb e_{31} 
	= \vb e_2 \vb e_3 \vb e_3 \vb e_1
	= \vb e_2 \vb e_1
	= - \vb e_{12}
	\,,
\\
	\vb e_{31} \vb e_{23} 
	= \vb e_3 \vb e_1 \vb e_2 \vb e_3
	= \vb e_1 \vb e_3 \vb e_3 \vb e_2
	= \vb e_1 \vb e_2
	= \vb e_{12}
	\,,
\\
	{\vb e_{31}}^2 
	= \vb e_3 \vb e_1 \vb e_3 \vb e_1 
	= - \vb e_3 \vb e_3 \vb e_1 \vb e_1 
	= - \mathbb 1
	\,,
\\
	{\vb e_{23}}^2 
	= \vb e_2 \vb e_3 \vb e_2 \vb e_3 
	= - \vb e_2 \vb e_2 \vb e_3 \vb e_3 
	= - \mathbb 1
	\,,
\end{gather}
o algebroje $A_0^-$
\begin{gather}
	\vb e_{23} \vb e_{31} 
	= \vb e_2 \vb e_3 \vb e_3 \vb e_1 
	= - \vb e_2 \vb e_1 
	= \vb e_1 \vb e_2 
	= \vb e_{12}
	\,,
\\
	\vb e_{31} \vb e_{23} 
	= \vb e_3 \vb e_1 \vb e_2 \vb e_3
	= \vb e_1 \vb e_3 \vb e_3 \vb e_2
	= - \vb e_1 \vb e_2
	= - \vb e_{12}
	\,,
\\
	{\vb e_{31}}^2 
	= \vb e_3 \vb e_1 \vb e_3 \vb e_1 
	= - \vb e_3 \vb e_3 \vb e_1 \vb e_1 
	= - \mathbb 1
	\,,
\\
	{\vb e_{23}}^2 
	= \vb e_2 \vb e_3 \vb e_2 \vb e_3 
	= - \vb e_2 \vb e_2 \vb e_3 \vb e_3 
	= - \mathbb 1
	\,.
\end{gather}
Abiems algebroms $\vb e_{31}$ ir $\vb e_{23}$ yra algebros generatoriai.
Jie tenkina sąryšius
\begin{equation}
	\begin{aligned}
		\vb e_{31} \vb e_{23} &= - \vb e_{23} \vb e_{31} \,,\\
		{\vb e_{31}}^2 &= - \mathbb1 \,,\\
		{\vb e_{23}}^2 &= - \mathbb1 \,.
	\end{aligned}
\end{equation}

Tegul $f:A^+_0 \to A^-_0$ yra tiesinis atvaizdavimas,
apibrėžiamas veikimu ant $A^+_0$ bazės
\begin{equation}
	\mathbb1 \mapsto \mathbb1
	\,,\quad
	\vb e_{23} \mapsto -\vb e_{23}
	\,,\quad
	\vb e_{31} \mapsto -\vb e_{31}
	\,,\quad
	\vb e_{12} \mapsto -\vb e_{12}
	\,.
\end{equation}
Akivaizdu kad $f$ yra bijekcija.
Įrodymui kad $f$ yra algebrų homomorfizmas užtenka parodyti,
kad $f$ nepakeičia generatorių daugybos sąryšių
\begin{equation}
	\begin{aligned}
		f(\vb e_{31}) f(\vb e_{23}) 
		=(-\vb e_{31})(-\vb e_{23}) 
		&= - \vb e_{23} \vb e_{31} 
		= - f(\vb e_{23}) f(\vb e_{31}) 
\,,\\
		f(\vb e_{31})^2 
		=(-\vb e_{31})^2 
		&= - \mathbb1
		= - f(\mathbb1)
\,,\\
		f(\vb e_{23})^2 
		=(-\vb e_{23})^2 
		&= - \mathbb1
		= - f(\mathbb1)
\,.
	\end{aligned}
\end{equation}
Atvaizdavimas $f$ yra algebrų izomorfizmas,
tad $A^+_0 \cong A^-_0$.
\end{proof}


\subsection{Grupės $\pinp$ ir $\pinm$}

Algebros $\mathup{Cl^\pm(V)}$ poaibis sudarytas iš elemtų $\vb v_1 \vb v_2
\dots \vb v_N$, kur visi $\vb v_i$ tenkina $(\vb v_i,\vb v_i) = 1$, kartu su
algebros daugyba sudaro atitinkamą grupę $\pinpm$. Grupės uždarumo sąvybė akivaizdžiai tenkinama, 
vienetinis elementas yra tas pats kaip ir algebros vienetinis elementas $\mathbb 1$,
o kaip gaunamas atvikštinis elementas siriasi grupeių $\pinp$ ir $\pinm$.

Grupėje $\pinp$ elemento $v$ atvirkštinis elementas yra tranponuotas $v^t$
\begin{multline}
	\vb v_1 \vb v_2 \dots \vb v_N (\vb v_1 \vb v_2 \dots \vb v_N )^t
	=
	\vb v_1 \vb v_2 \dots \vb v_N \vb v_N \vb v_{N-1} \dots \vb v_1
	=
	(\vb v_N,\vb v_N)
	(\vb v_{N-1},\vb v_{N-1})
	\dots
	(\vb v_1,\vb v_1)
	\mathbb 1
	=
	\mathbb 1
\end{multline}
Grupėje $\pinm$ atvirkštinis elementas yra $\alpha(v^t)$
\begin{multline}
	\vb v_1 \vb v_2 \dots \vb v_N \alpha\qty(\vb v_1 \vb v_2 \dots \vb v_N )^t
	=
	\vb v_1 \vb v_2 \dots \vb v_N \alpha(\vb v_N) \alpha(\vb v_{N-1}) \dots \alpha(\vb v_1)
	\\=
	\vb v_1 \vb v_2 \dots \vb v_N (-\vb v_N) (-\vb v_{N-1}) \dots (-\vb v_1)
	=
	(\vb v_N,\vb v_N)
	(\vb v_{N-1},\vb v_{N-1})
	\dots
	(\vb v_1,\vb v_1)
	\mathbb 1
	=
	\mathbb 1
\end{multline}

Grupė $\spin$ galima gauti dviem būdais: kaip $\pinp$ pogrupį sudarytą
iš lyginių vektorių sandaugų arba kaip $\pinm$ pogrupį gautą taip pat.
Abi konstrukcijos yra viena kitai izomofiškos, nes minusas iš sandaugos
$\pinm$ pasirodys lyginį kartą sakičių ir rezultatas bus toks pat kaip
$\pinp$ pogrupyje.

\subsection{$\spin$ ir $\pinpm$ veikimas ant $\set R^3$}

Visos aptartos grupės veikia vektorius iš $\set R^3$, juos pirma patalpinant į
$\mathup{Cl^\pm(V)}$ ir tada su algebros sandauga
\begin{equation}
	\rho(v).x = \begin{cases}
		vxv^t		& v \in \pinp \\
		\alpha(v)xv^t	& v \in \pinm 
	\end{cases}
\end{equation}






\subsection{Clifordo algebros}


Toliau nagrinėsime nagrinėsime dvi Clifordo algebras $\mathup{Cl}^+(3)$ ir
$\mathup{Cl}^-(3)$. Tai yra vectorinės erdvės su papildoma bitiesine daugybos
operacija, kuri leidžia sudauginti du algebros vektorius ir gauti naują
vektorių algebroje. Šias algebras apibrėšime pirma paimdami daug bendresnę
tenzorinę algebrą ir tada jai sutekdami norimą struktūrą.

\begin{definition}
	Vienetinė tenzorinė algebra $T(\set R^3)$ yra tiesioginė suma
	\begin{equation}
		T(\mathbb{R}^3) 
		=
		\bigoplus_{k=0}^{\infty} T^k
		=
		\mathbb{R} 
		\oplus \set R^3 
		\oplus (\set R^3 \otimes \set R^3) 
		\oplus (\set R^3 \otimes \set R^3 \otimes \set R^3) 
		\oplus \dots
		\,,
	\end{equation}
	$T^k$ čia yra vektorinės erdvės $T^0 = \set R$\: ,\: $T^1=\set R^3$\: ,\: $T^2=\set R^3 \otimes \set R^3$\: ,\:
	$T^3=\set R^3 \otimes \set R^3 \otimes \set R^3$
	ir taip toliau. Algebros daugyba yra duodama vektorių tenzorinės sandaugos.
	Vienetinis tenzorinės algebros elementas yra žymimas $\mathbb 1$\:.
\end{definition}
Kitaip paskius tai yra algebra, kurios vektoriai yra ivairių ilgiu tenzorinių
sandaugos $\vb v_1 \otimes \vb v_2 \otimes \dots \otimes \vb v_n$\:,\: $\vb v_i \in \set R^3$
ir tokių sandaugų sumos.
\begin{example}
	Vektorius $3\cdot\mathbb 1 + \vb v$ priklauso $T(\set R^3)$ ir yra suma vektorių iš $T^0$ ir $T^1$.
	Vektorius $2\vb w + \vb v \otimes \vb w$ yra suma vektorių iš $T^1$ ir $T^2$.
	Jų sandauga tenzorinėje algebroje yra
	\begin{equation}
		\begin{split}
			\qty(3\cdot\mathbb 1 + \vb v)
			\otimes
			\qty(2\vb w + \vb v \otimes \vb w)
			=&
			3\cdot\mathbb 1 \otimes 2\vb w
			+ 3\cdot\mathbb 1 \otimes \vb v \otimes \vb w
			+ \vb v \otimes 2\vb w
			+ \vb v \otimes \vb v \otimes \vb w
			\\=&
			3\cdot2\qty(\mathbb 1 \otimes \vb w)
			+ 3\qty(\mathbb 1 \otimes \vb v \otimes \vb w)
			+ 2(\vb v \otimes \vb w)
			+ \qty(\vb v \otimes \vb v \otimes \vb w)
			\\=&
			6 \vb w
			+ 3\qty(\vb v \otimes \vb w)
			+ 2(\vb v \otimes \vb w)
			+ \qty(\vb v \otimes \vb v \otimes \vb w)
			\\=&
			6 \vb w
			+ 5\qty(\vb v \otimes \vb w)
			+ \qty(\vb v \otimes \vb v \otimes \vb w)
			\,,
		\end{split}
	\end{equation}
	čia pirmoje eilutėje pasinaudojome tuo kad tenzorinės sandauga
	distributyvi; antroje eilutėje naudojames tuo kad tenzorinė sandauga
	tiesiška kiekvieno dauginamo vektoriaus atžvilgiu, kad iškeltųme 
	skaliarinius daugiklius; trečioje eilutėje sudedame tiesiškai
	priklausomus vektorius.
	Gautas vektorius yra suma vektorių iš $T^1$, $T^2$ ir $T^3$.
\end{example}


Dabar $T(\set R^3)$ suteiksime papildomą strutūrą, kad algebroje
$\mathup{Cl}^+(3)$ būtų tenkinamas sąryšis
\begin{equation}
	\vb v \otimes \vb w
	+
	\vb w \otimes \vb v
	=
	2(\vb v , \vb w) \mathbb 1
	\quad ,
	\label{eq:cl+}
\end{equation}
o algebroje $\mathup{Cl}^+(3)$ tenkinamas sąryšis
\begin{equation}
	\vb v \otimes \vb w
	+
	\vb w \otimes \vb v
	=
	- 2(\vb v , \vb w) \mathbb 1
	\quad .
	\label{eq:cl-}
\end{equation}
Šiuos sąryšius ,,užkoduojame'' į aibes $S^\pm \subset T(\set R^3)$
\begin{equation}
	\begin{split}
		S^+ = \qty{
			\vb v \otimes \vb w
			+
			\vb w \otimes \vb v
			-
			2(\vb v , \vb w) \mathbb 1
			\mid
			\vb v, \vb w \in \set R^3
		}
		\\
		S^- = \qty{
			\vb v \otimes \vb w
			+
			\vb w \otimes \vb v
			+
			2(\vb v , \vb w) \mathbb 1
			\mid
			\vb v, \vb w \in \set R^3
		}\,.
	\end{split}
\end{equation}
Užrašome tenzorinės algebros poalgebrius $I^\pm \subset T(\set R^3)$,
gauamus iš visų vektorių, kurie turi bent vieną daugiklį iš $S^\pm$
\begin{equation}
	I^+ = \qty{
		\sum_{\substack{
			s \in S^\pm \\ 
			a,b \in \mathup{T(V)} 
		}}
		a \otimes s \otimes b
	}
	\quad .
\end{equation}
Šie poalbebriai papildomai tenkina savybę, kad betkokiem $a \in T(\set R^3)$
ir $d \in I^\pm$, $a\otimes d, d \otimes a \in I^\pm$. Poalgebris
tenkinantis šia savybę yra vadinamas algebros idealu.
Literatūroje šią konstrukciją yra įprasta pateikti sakiniu:
,,$I^\pm$ yra algebros $T(\set R^3)$ idealas genaruojams 
sąryšio (\ref{eq:cl+}) ar atitinkamai (\ref{eq:cl-})''.

\begin{definition}
	Algebra $\mathup{Cl}^+(3)$ yra faktorinė algebra $T(\set R^3)/I^+$
	ir algebra $\mathup{Cl}^-(3)$ yra faktorinė algebra $T(\set R^3)/I^-$.
\end{definition}
Faktorinės algebros elementai yra šoninės klasės sudarytos iš elementų,
kuriuos galime gauti vieną iš kito pridedant vektorių iš idealo.
\begin{example}
	$2\vb{(v,w)} \mathbb 1 - \vb v \otimes \vb w + I^+$ yra šoninė klasė,
	kuri yra vektorius algebroje $\mathup{Cl}^+(3)$. 
	$2\vb{(v,w)} \mathbb 1 - \vb v \otimes \vb w$ 
	yra šoninės klasės atstovas, o ,,$+ I^+$'' žymi
	kad šiai klasei taip pat prikauso visi vectoriai gaunami prie
	klasės atstovo pridėjus vektorių iš $I^+$.
	Tai mums leidžia atlikti skaičiavimus kaip
	\begin{equation}
		\vb{(v,w)} \mathbb 1 - \vb v \otimes \vb w + I^+
		= 
		\vb{(v,w)} \mathbb 1 - \vb v \otimes \vb w 
		+ \qty(
			\vb v \otimes \vb w + \vb w \otimes \vb v 
			- 2\vb{(v,w)}\mathbb 1
		) 
		+ I^+
		=
		\vb w \otimes \vb v + I^+
		\,,
	\end{equation}
	čia pasinaudojome tuo kad 
	$\qty( \vb v \otimes \vb w + \vb w \otimes \vb v - 2\vb{(v,w)}\mathbb 1) \in I^+$.
	Pavyko parodyti, kad $2\vb{(v,w)} \mathbb 1 - \vb v \otimes \vb w$ 
	ir $\vb w \otimes \vb v$ atstovaują tą pačią klasę.
\end{example}

\begin{example}
	Tegul $a,b \in T(\set R^3)$ ir
	$
		a \otimes \vb{(v \otimes w + v \otimes w)} \otimes b + I^-
	$
	yra šoninė klasė, kuri yra vektorius iš $\mathup{Cl}^-(3)$.
	Šią išraišką supapranstiname atlidami veiksmus
	\begin{equation}
		\begin{split}
			a \otimes \vb{(v \otimes w + v \otimes w)} \otimes b + I^-
			=&
			a \otimes \vb{(v \otimes w + v \otimes w)} \otimes b 
			- a \otimes \qty(\vb{v \otimes w + v \otimes w}+2\vb{(v,w)} \mathbb 1)  \otimes b 
			+ I^-
			\\=&
			a \otimes \qty(
				\vb{(v \otimes w + v \otimes w)}
				-\qty(\vb{v \otimes w + v \otimes w}
				-2\vb{(v,w)} \mathbb 1) 
			)\otimes b 
			+ I^-
			\\=&
			a \otimes \qty(
				-2\vb{(v,w)} \mathbb 1 
			)\otimes b 
			+ I^-
			\\=&
			-2\vb{(v,w)} \qty(a \otimes \mathbb 1 \otimes b)
			+ I^-
			\\=&
			-2\vb{(v,w)} \qty(a \otimes b)
			+ I^-
			\,.
		\end{split}
	\end{equation}
	Parodėme, kad sudarius faktorinę algebrą $\mathbb{Cl}^-(3)$,
	galime jos viduje taikyti sąryšį (\ref{eq:cl-})
	manipuliuojant šoninės klasės atstovus.
\end{example}
\begin{example}
	Tegul $a$ ir $b$ yra šoninių klasių atstovai iš
	$\mathup{Cl}^\pm(3)$.
	Jų sandauga algebroje yra
	\begin{equation}
		(a+I^\pm)\otimes(b+I^\pm)
		=
		a \otimes b
		+ a \otimes I^\pm
		+ I^\pm \otimes b
		+ I^\pm \otimes I^\pm
		=
		a \otimes b + I^\pm
		\,.
	\end{equation}
	Sumos nariai $
		a \otimes I^\pm
		= I^\pm \otimes b
		= I^\pm \otimes I^\pm
		= I^\pm
	$, nes $I^\pm$ yra algebros $T(\set R^3)$ idealas.
	Tai
	\begin{equation}
		(a+I^\pm)\otimes(b+I^\pm)
		=
		a \otimes b + I^\pm
		\,.
	\end{equation}
\end{example}
Toliau priimame sutrumpintą žymėjimą, nes išreikštai
su šoninėmis klasėmis nebedirbsime. Užrašysime tik
kalsės atstovą, o kad išraiška skirtūsi nuo 
vektorių tenzorinės sandaugos vietoje ženklo ,,$\otimes$''
tiesiog rašysime vektorius vieną šalia kito.
Vietoje
\begin{equation}
	a \otimes \vb{(v \otimes w + v \otimes w)} \otimes b + I^-
	=
	-2\vb{(v,w)} \qty(a \otimes b) + I^-
\end{equation}
rašysime
\begin{equation}
	a \vb{(v w + v w)} b
	=
	-2\vb{(v,w)} a b
	\,.
\end{equation}

Štai kelios svarbios tapatybės
\begin{preposition}
	\begin{equation}
		\vb v \vb v 
		= \frac{1}{2}(\vb v \vb v + \vb v \vb v) 
		= \pm(\vb v,\vb v) \mathbb 1
		\,.
	\end{equation}
\end{preposition}

\begin{preposition}
	Algebroje $\mathup{Cl}^+(3)$ 
	\begin{equation}
		-\vb v \vb w \vb v 
		= -\qty( 2(\vb v,\vb w)\mathbb 1 - \vb{wv} )\vb v 
		= -2\vb{(v,w)} \vb v + \vb{wvv} 
		= -2\vb{(v,w)} \vb v + \vb{(v,v)}\vb w 
		\,.
	\end{equation}
\end{preposition}

\begin{preposition}
	Algebroje $\mathup{Cl}^-(3)$ 
	\begin{equation}
		\vb v \vb w \vb v 
		= \qty( - 2(\vb v,\vb w)\mathbb 1 - \vb{wv} )\vb v 
		= -2\vb{(v,w)} \vb v - \vb{wvv} 
		=  -2\vb{(v,w)} \vb v + \vb{(v,v)}\vb w 
		\,.
	\end{equation}
\end{preposition}

Atkreipiame dėmesį, kad kai $(v,v)=1$,
$\vb{(v,w)v}$ yra vektoriaus $w$ projekcija į $v$
ir $-2\v{(v,w)v} + \vb w$ yra vektoriaus atspindys
per plokštumą statmeną vektoriui $v$.

Tegul $\qty{\vb e_1,\vb e_2,\vb e_3}$ yra ortonormali $\set R^3$ bazė. Bazinius vektorius galime perketi
į algebra $\mathup{Cl^\pm(V)}$, kur jie turi tenkinti
\begin{equation}
	\begin{split}
		\vb e_i \vb e_j &= - \vb e_j \vb e_i \ , \ i \neq j \\
		\vb e_i \vb e_i &= \pm \mathbb 1
		\quad .
	\end{split}
\end{equation}
Su šių sąryšių pagalba galime bet kokį $\mathup{Cl^\pm(V)}$ elementą užrašyti, kaip
tiesinę kombinaciją aštuonių elementų:
\begin{equation}
		\mathbb 1 ,\,
		\vb e_1 ,\, \vb e_2 ,\, \vb e_3 ,\, 
		\vb e_2 \vb e_3 ,\, \vb e_1 \vb e_3 ,\, \vb e_1 \vb e_2 ,\,
		\vb e_1 \vb e_2 \vb e_3 ,
\end{equation}
nes su pirma taisykle betkurį tenzorinės algebros bazės elementą galima išrikiuoti taip kad
ortų indeksai būtų surašyti didėjimo tvarka ir su antra taisykle pasikartojantys ortai panaikinami.
Tai reiškia $\mathup{Cl^\pm(V)}$ yra aštuonmatės algebros. Algebros elemento norma $\abs{v}$ yra
abibrėžiama, taip pat kaip vektoriaus ilgis
\begin{equation}
	\abs{v} = \sqrt{ \sum_{I \subset {1,2,3}} {v_I}^2 }
	\quad ,
\end{equation}
čia summa eina per visus indeksų poaibius, kurie atitinka bazinius algebgros elementus.
Bendresniu atvėju $\mathup{dim}\qty[\mathup{Cl}^\pm\qty(\set R^n)] = 2^n$.

Apibrėžiame "transponuotą elementą" apibėždami kaip transponavimas veikia ant algebros bazės
\begin{equation}
	\qty(\vb e_{i_1} \vb e_{i_2} \dots \vb e_{i_k} )^t
	=
	\vb e_{i_k} \vb e_{i_{k-1}} \dots \vb e_{i_1}
	=
	(-1)^{k-1} \vb e_{i_1} \vb e_{i_2} \dots \vb e_{i_k} 
	\quad .
\end{equation}
Transpozicija yra algebros anti-homomorfizmas, tai reiškia, kad betkuriems algebros elementams $v$ ir $w$ galioja
\begin{equation}
	(v w)^t = w^t v^t
	\quad .
\end{equation}


Taip pat veikimu ant bazės apibrėžiame algebros homomorfizmą
\begin{equation}
	\alpha( \vb e_{i_1} \vb e_{i_2} \dots \vb e_{i_k} )
	= 
	(-1)^{k} \vb e_{i_1} \vb e_{i_2} \dots \vb e_{i_k} 
	\quad ,
\end{equation}
jis algebros elemento komponentes prie lyginio ortų skaičiau siunčia į savo minusą, o
komponentčių prie nelyginio skaičiaus nepakeičia.













\subsection{Grupės $\pinp$ ir $\pinm$}

\begin{definition}
	$\pinpm$ yra grupės genruojamos vienetiniu vektorių $\vb v \in R^3$,
	$(\vb v,\vb v) = 1$, juos dauginat algebroje $\Clpm$.
	Gupės elemnetai yra vienanariai $\vb v_1 \vb v_2 \dots \vb v_n$
	ir gupės daugyba sutampa su algebros daugyba.
\end{definition}

\begin{definition}
	Grupė $\spin$ yra pogrupis $\pinp$,
	sudarytas iš visų lyginių $\pinp$ elementų.
\end{definition}

\begin{preposition}
	Jie ir tik jei $v$ priklauso $\pinp$,
	$vv^t=\mathbb1$.
	Jei ir tik jei $v$ priklauso $\pinm$,
	$v\alpha(v^t)=\mathbb1$.
\end{preposition}

\begin{example}
	Patikrinkime, kad $\pinm$ tikrai tenkina grupės aksiomas.
	\begin{itemize}
		\item Asociatyvumas yra paveldimas iš algebros
			$\Clm$,
			kuri yra asociatyvi;
		\item tapatybės elementas yra $\mathbb 1$;
		\item atvikštinis elementas yra $v^{-1}= \alpha(v^t)$;
		\item sandaugos uždarumą patikriname skaičiavimu
	\end{itemize}
	\begin{align*}
		\alpha((vw)^t) vw &= \alpha(w^t v^t) vw && \text{antihomomorfizmo savybė}
		\\
		&= \alpha(w^t) \alpha(v^t) vw && \text{homomorfizmo savybė}
		\\
		&= \alpha(w^t) \mathbb 1 w && \text{iš $\pinm$ apibrėžimo}
		\\
		&= \mathbb 1
		\,.
	\end{align*}
\end{example}

Grupėms $\pinpm$ prikkauso visi vienetiniai vektoriai
perkelti iš $\set R^3$ į $\mathrm{Cl}^\pm(3)$.
Jie yra $\pinpm$ generatoriai,
tad visus $\pinpm$ elementus galima realizuoti,
kaip vienetinių vectorių sandaugą algebroje
$\mathrm{Cl}^\pm(3)$. $\set Z_2$ gradavimas išlika
ir ant $\pinpm$.
Elementai $\vb v$ ir $\vb{vwu}$ yra nelyginiai,
o elementai $\mathbb 1$ ir $\vb{vw}$ yra lyginiai.
Dauginant tarpusavyje lyginius elementus,
gausime tik lyginius elementus,
todėl jie sudaro $\pinpm$ pogrupį.


\begin{preposition}
	Grupių $\pinp$ ir $\pinm$ lyginiai pogrupiai
	yra izomorphiški, kaip grupės.
\end{preposition}
\begin{proof}
	$\pinpm$ lyginiai porgupiai kaip aibės prikauso 
	atitinkamiems $\Clpm$ lyginiams poalgiabriams $A_0^\pm$.
	Poalgebriai $A_0^+ \cong A_0^-$ yra izomorfiški kaip algebros (teiginys \ref{prep:even_iso}),
	tegul $f: A_0^+ \to A_0^-$ žymi izomorphizmą tarp jų.
	Jei $\Clp$ lyginis elementams $v$ priklauso $\pinp$,
	tai $f(v)$ priklauso $\pinm$
	\begin{equation}
		\begin{aligned}
			v^t *_+ v &= \mathbb 1_+
			\\
			f(v^t *_+ v ) &= f\qty(\mathbb 1_+)
			\\
			f( v^t) *_- f(v) &= \mathbb 1_-
			&& \text{izomorfizmo savybė}\\
			{f(v)}^t *_- f(v) &= \mathbb 1_-
			&& \text{transponavimas gerbia isomorfizmą}\\
			\alpha\qty({f(v)}^t) *_- f(v) &= \mathbb 1_-
			&& f \text{ kodomenas yra lyginiai elementai}
			\,,
		\end{aligned}
	\end{equation}
	čia $*_+$ ir $*_-$ yra algebros daugybos atitinkamose algebrose,
	o $\mathbb 1_+$ ir $\mathbb 1_-$ yra 
	atitinkamų algebrų vienetiniai elementai.
	Atitinkamai gauname rezultatą, kad jei $v$ prikauso $\pinp$,
	tai $f^{-1}(v)$ priklauso $\pinm$.
	Apribojus domeną ir kodomeną funksija $f$ yra
	grupių izomorfizmas, nes grupės sandauga
	sutampa su algebros sandauga.
\end{proof}

Tiek $\pinp$, tiek $\pinm$ lyginius pogrupius vadinsime vienodai - $\spin$,
nes jie yra izomofiški.

\begin{preposition}
	Visi $\spin$ elementai yra formos
	\begin{equation}
		a_0 \mathbb1 
		+ a_1 \vb e_{23}
		+ a_2 \vb e_{31}
		+ a_3 \vb e_{12}
	\end{equation}
	kur $a_i \in \set R$ ir
	\begin{equation}
		{a_0}^2
		+{a_1}^2
		+{a_2}^2
		+{a_3}^2
		= 1
	\end{equation}
\end{preposition}
\begin{proof}
	Bendra lyginio $\Clp$ vektoriaus forma yra
	\begin{equation}
		v = 
		  a_0 \mathbb1 
		+ a_1 \vb e_{23}
		+ a_2 \vb e_{31}
		+ a_3 \vb e_{12}
	\end{equation}
	Kad elementas $v$ priklausytų $\spin$, jis turi tenkinti
	$vv^t = \mathbb1$.
	Įsisteatę $v$ bendra formą gauname
	\begin{equation}
		\begin{aligned}
			(a_0 \mathbb1 
			+ a_1 \vb e_{23}
			+ a_2 \vb e_{31}
			+ a_3 \vb e_{12})
			(a_0 \mathbb1 
			- a_1 \vb e_{23}
			- a_2 \vb e_{31}
			- a_3 \vb e_{12})
			&= \mathbb1
			\\
			({a_0}^2
			+{a_1}^2
			+{a_2}^2
			+{a_3}^2)
			\mathbb1
			&= \mathbb1
		\end{aligned}
	\end{equation}
	Rezultatas seka iš fakto kad ${\vb e_{ij}}^2 = -\mathbb1$ 
	ir $\vb e_{ij} \vb e_{kl} = - \vb e_{kl} \vb e_{ij}$.
\end{proof}



\subsection{$\spin$ ir $\pinpm$ veikimas ant $\vecSpace R 3$}

Visos aptartos grupės veikia vektorius iš $\set R^3$, juos pirma patalpinant į
$\mathup{Cl^\pm(V)}$ ir tada veikiant su algebros sandauga
\begin{equation}
	\rho(v).\vb w := \alpha(v) \vb w v^{-1} 
	= \begin{cases}
		\alpha(v) \vb w v^t		& v \in \pinp \\
		v \vb w v^t	& v \in \pinm 
	\end{cases}
\end{equation}

Iškart galima pamatyti kad $\pinpm$ elementai $\vb v$ ir $-\vb v$ veikia vienodai,
nes daugiklis $-1$ du kart iškeliamas prieš algebros sandaugą.
Teiginiai \ref{prep:R_action+} ir \ref{prep:R_action-} parodo kad
vienetinis vektorius $\vb n \in \set R^3$ patalpinti į $\pinpm$,
veikia ant $\vb w \in \set R^3$,
kaip atspindys per plokštumą statmeną vektoriui $\vb n$.
Tada lyginiai elementai $\vb v \vb w$ ant $\set R^3$ veikia,
kaip atspindys per vieną plokštumą tada atspindys per kitą
plokštumą. Tai yra teisingas pasukimas.
Kitas elementas kurio veikimą svarbu patikrinti yra
$\vb e_{123}$, gauname $\rho(\vb e_123).\vb e_i = - e_i$,
tai $e_{123}$ ant $\set R^3$ veikia kaip inversija.





\begin{theorem}[{\cite[Theorem~2.9]{lawson1989spin}}, atvėjis kai $V=\set R^3$ ir $q(v)=\pm(v,v)$]
	Egzistuoja truposios tiksliosios sekos (ang. \emph{short exact sequence})
	\begin{equation}
		\begin{tikzcd}
			1 \arrow[r] &\{1,-1\} \arrow[r]	&\pinp \arrow[r, "\rho"] &\Orth3 \arrow[r] &1
		\end{tikzcd}
	\end{equation}
	\begin{equation}
		\begin{tikzcd}
			1 \arrow[r] &\{1,-1\} \arrow[r]	&\pinm \arrow[r, "\rho"] &\Orth3 \arrow[r] &1
		\end{tikzcd}
	\end{equation}
	\begin{equation}
		\begin{tikzcd}
			1 \arrow[r] &\{1,-1\} \arrow[r]	&\spin \arrow[r, "\rho"] &\SO3 \arrow[r] &1
		\end{tikzcd}
	\end{equation}
	Čia kiekviena rodyklė yra grupių homorfizmas ir kiekvieno homomorfizmo vaizdas
	yra branduolys homomorfizmo einančio po jo sekoje.
\end{theorem}




\subsection{$\spin$ ir $\pinpm$ matriciniai įvaizdžiai}

Dėl patogumo dirbant su grupėmis ir algebromis apibrėšime
nagrinėjamų grupių ir algebrų ištikimus matricinius įvaizdžius.
Šiuos matricinius įvaizdžius užrašysime su Pauli matricų 
ir vienetinės $2\times2$ matricos pagalba
\begin{equation}
	I := \mqty(\imat{2})
	\,,\qquad
	\sigma_1 := \mqty(\pmat{1}) 
	\,,\qquad
	\sigma_2 := \mqty(\pmat{2}) 
	\,,\qquad
	\sigma_3 := \mqty(\pmat{3}) 
	\,.
\end{equation}
Šios matricos tenkina
\begin{align}
	\sigma_i \sigma_i &= I \,,
	&&i\in{1,2,3}\,;
	\\
	\sigma_i \sigma_j &= - \sigma_j \sigma_i\,, 
	&&i\neq j \,;
	\\
	\sigma_i \sigma_j &= i \sigma_k \,,
	&&\text{$i,j,k$ yra indeksų $1,2,3$ cikliniai perstatymai}.
\end{align}


\begin{preposition}
	\label{prep:pauli_rep+}
	Veikimu ant $\Clp$ \emph{lyginio} poalgebrio $A_0^+$ bazės apibrėžtas atvaizdavimas
	\begin{align*}
		\mathbb 1	\mapsto	I		\qquad
		\vb e_{23} 	\mapsto	i\sigma_1	\qquad
		\vb e_{31}	\mapsto	i\sigma_2	\qquad
		\vb e_{12}	\mapsto	i\sigma_3	
	\end{align*}
	yra algebros $\Clp$ \emph{lyginio} poalgebrio ištikimas matricinis įvazdis.
\end{preposition}
\begin{proof}
	Užtenka pratikrinti kad pasirinkti algebros generatoriai $\vb e_{23}$ ir $\vb e_{31}$
	tapusavyje tenkina tokius pat sandaugos sąryšius kaip jų įvazdzdis.
	\begin{align}
		{\vb e_{23}}^2 = {\vb e_{31}}^2 &= - \mathbb1
		\\
		\vb e_{23} \vb e_{31} &= - \vb e_{12}
		\\
		\vb e_{31} \vb e_{23} &= \vb e_{12}
	\end{align}
	\begin{align}
		\qty(i\sigma_1)^2 = \qty(i\sigma_2)^2 = - {\sigma_2}^2 &= - I
		\\
		i\sigma_1 i\sigma_2 = - \sigma_1 \sigma_2 &= - i\sigma_3
		\\
		i\sigma_2 i\sigma_1 = - \sigma_2 \sigma_1 = \sigma_1 \sigma_2 &= i\sigma_3
	\end{align}
\end{proof}


\begin{preposition}
	\label{prep:pauli_rep-}
	Veikimu ant $\Clm$ \emph{lyginio} poalgebrio $A_0^-$ bazės apibrėžtas atvaizdavimas
	\begin{align*}
		\mathbb 1	\mapsto	I		\qquad
		\vb e_{23} 	\mapsto	-i\sigma_1	\qquad
		\vb e_{31}	\mapsto	-i\sigma_2	\qquad
		\vb e_{12}	\mapsto	-i\sigma_3	
	\end{align*}
	yra algebros $\Clm$ \emph{lyginio} poalgebrio ištikimas matricinis įvazdis.
\end{preposition}
\begin{proof}
	Analogiškai kaip teiginyje \ref{prep:pauli_rep+}
\end{proof}

\begin{preposition} 
	Atvizdis iš teiginio \ref{prep:pauli_rep+} siunčia $\spin$ į matricinė grupė $\mathup{SU}(2)$.
	Tai yra $\spin$ ir $\mathup{SU}(2)$ yra izomorfiškos grupės.
\end{preposition}
\begin{proof}
	Bendros formos $\spin$ elementas
	$ a_0 \mathbb1 + a_1 \vb e_{23} + a_2 \vb e_{31} + a_3 \vb e_{12}$,
	$ {a_0}^2 +{a_1}^2 +{a_2}^2 +{a_3}^2 = 1 $,
	yra atvaizduojamas į maricą
	\begin{equation}
		a_0 I + a_1 i\sigma_1 + a_2 i\sigma_2 + a_3 i\sigma_3
		=
		\mqty(
			  a_0 + ia_3	&& a_2 + ia_1 \\
			- a_2 + ia_1	&& a_0 - ia_3 
		)
		=
		\mqty(
			z_1	&& -z_2^* \\
			z_2	&& z_1^* 
		)
		\,,
	\end{equation}
	čia pasižymėjome kompleksinius skaičius 
	$z_1 = a_0 + i a_3$ ir $z_2 = - a_2 + ia_1$. 
	Sąlyga $ {a_0}^2 +{a_1}^2 +{a_2}^2 +{a_3}^2 = 1 $
	tada tampa $ \abs{z_1}^2 + \abs{z_2}^2 = 1$.
	Kas reiškia kad $
		\smqty(
			z_1	&& -z_2^* \\
			z_2	&& z_1^* 
		)
	$ yra benda forma matricos, kuri priklauso $\mathup{SU}(2)$.
\end{proof}

Gerai žinoma jog $\mathup{SU}(2)$ elmentas atitinkantis pasukimą
kampu $\theta$ aplink vienetinį vektoriu $\vb n$ yra
\begin{equation}
	e^{\frac{i\theta}{2} (\vb n,\symbf\sigma)}
	\,,
\end{equation}
čia $(\vb n,\symbf\sigma) = n_1 \sigma_1 + n_2 \sigma_2 + n_3 \sigma_3 $
ir kvadratinės matricos eksponetė yra apibrėžta kaip 
\begin{equation}
	e^X := \sum_{n=0}^\infty \frac{X^n}{n!}
	\,.
\end{equation}

Algebrų $\Clpm$ įvaizdžiams užrašyti pasinaudosime gama maricomis,
kurios chiralėje bazėja yra apibrežiamos kaip
\begin{equation}
	\gamma^0 := \mqty(\dmat{I,I}) 			\,,\quad
	\gamma^1 := \mqty(\admat{\sigma_1,-\sigma_1})	\,,\quad
	\gamma^2 := \mqty(\admat{\sigma_2,-\sigma_2})	\,,\quad
	\gamma^3 := \mqty(\admat{\sigma_3,-\sigma_3})	\,.
\end{equation}
Apsiribosime tiktai maricomis $\gamma^1$, $\gamma^2$ ir $\gamma^3$.
Jos tenkina sąryšius
\begin{equation}
		{\gamma^i}^2 = -\mqty(\dmat{I,I})
		\,,\quad
		\gamma^i \gamma^j = - \gamma^j \gamma^i
		\,.
\end{equation}


\begin{preposition}
	Veikimu ant $\Clm$ bazės apibrėžtas atvaizdavimas
	\begin{align*}
		\mathbb 1	\mapsto				&\mqty(\dmat{I,I})			&\quad	\vb e_{123}	\mapsto	\gamma^1\gamma^2\gamma^3 =	&\mqty(\admat{-iI,iI})			\\
		\vb e_{23}	\mapsto	\gamma^2 \gamma^3 =	&\mqty(\dmat{-i\sigma_1,-i\sigma_1})	&\quad	\vb e_1 	\mapsto	\gamma^1 =			&\mqty(\admat{\sigma_1,-\sigma_1})	\\
		\vb e_{31}	\mapsto	\gamma^3 \gamma^1 =	&\mqty(\dmat{-i\sigma_2,-i\sigma_2})	&\quad	\vb e_2 	\mapsto	\gamma^2 =			&\mqty(\admat{\sigma_2,-\sigma_2})	\\
		\vb e_{12}	\mapsto	\gamma^1 \gamma^2 =	&\mqty(\dmat{-i\sigma_3,-i\sigma_3})	&\quad	\vb e_3 	\mapsto	\gamma^3 =			&\mqty(\admat{\sigma_3,-\sigma_3})	\\
	\end{align*}
	yra algebros $\Clm$ ištikimas matricinis įvazdis.
\end{preposition}
\begin{proof}
	Gama matricos tenkina tokius pat daugybos sąryšius,
	kaip algebros $\Clm$ generatoriai $\vb e_1$, $\vb e_2$ ir $\vb e_3$.
	\begin{align}
		{\vb e_i}^2 &= -\mathbb1		&\quad	{\gamma^i}^2 &= - \mqty(\dmat{I,I})\\
		\vb e_i \vb e_j &= -\vb e_j \vb e_i	&\quad	\gamma^i\gamma^j &= - \gamma^j\gamma^i
	\end{align}
	Nesunku pamatyti, kad bazinių algebros vectorių atvaizdžiai
	yra tiesiškai neprikausomos matricos.
	Iš to seką kad $\gamma^1$, $\gamma^2$, $\gamma^3$
	ir $\vb e_1$, $\vb e_2$, $\vb e_3$
	generuoja izomorfiškas algebras,
	tai yra turime ištikimą matricinį įvaizdį.
\end{proof}

Verta pabrėžti, kad lyginiai bazės elementai atvaizduojami į blok-diagonalias matricas,
kur kiekvienas blokas yra kopija įvaizdžio iš teiginio \ref{prep:pauli_rep-}.
Tai lyginį $\pinm$ elementą, kuris atinitnka sukimą apie $\vb n$ kampu $\theta$,
gauname kaip
\begin{equation}
		\mqty(\dmat{
			e^{-\frac{i\theta}{2}(\vb n,\symbf\sigma)},
			e^{-\frac{i\theta}{2}(\vb n,\symbf\sigma)}
		})
		=
		\exp(\frac{\theta}{2}(
			 n_1\gamma^2\gamma^3
			+n_2\gamma^3\gamma^1
			+n_3\gamma^1\gamma^2
		))
		\,.
\end{equation}
Visi nelyginiai elementai gali būti gauti,
kaip lyginis elementas padaugintas iš $\gamma^1\gamma^2\gamma^3 = \smqty(&-iI\\iI&)$,
tai nelyginio elemeto bendra forma yra
\begin{equation}
		\mqty(\admat{
			-ie^{-\frac{i\theta}{2}(\vb n, \symbf\sigma)},
			 ie^{-\frac{i\theta}{2}(\vb n, \symbf\sigma)}
		})
		=
		\gamma^1\gamma^2\gamma^3 \exp(\frac{\theta}{2}(
			 n_1\gamma^2\gamma^3
			+n_2\gamma^3\gamma^1
			+n_3\gamma^1\gamma^2
		))
		\,.
\end{equation}


\begin{preposition}
	Veikimu ant $\Clp$ bazės apibrėžtas atvaizdavimas
	\begin{align*}
		\mathbb 1	\mapsto				&\mqty(\dmat{I,I})			&\quad	\vb e_{123}	\mapsto	-i\gamma^1\gamma^2\gamma^3 =	&\mqty(\admat{I,-I})			\\
		\vb e_{23}	\mapsto	-\gamma^2 \gamma^3 =	&\mqty(\dmat{i\sigma_1,i\sigma_1})	&\quad	\vb e_1 	\mapsto	i\gamma^1 =			&\mqty(\admat{i\sigma_1,-i\sigma_1})	\\
		\vb e_{31}	\mapsto	-\gamma^3 \gamma^1 =	&\mqty(\dmat{i\sigma_2,i\sigma_2})	&\quad	\vb e_2 	\mapsto	i\gamma^2 =			&\mqty(\admat{i\sigma_2,-i\sigma_2})	\\
		\vb e_{12}	\mapsto	-\gamma^1 \gamma^2 =	&\mqty(\dmat{i\sigma_3,i\sigma_3})	&\quad	\vb e_3 	\mapsto	i\gamma^3 =			&\mqty(\admat{i\sigma_3,-i\sigma_3})	\\
	\end{align*}
	yra algebros $\Clp$ ištikimas matricinis įvazdis.
\end{preposition}
\begin{proof}
	Patikriname ar įvaizdis gerbia generatorių $\vb e_1$, $\vb e_2$ ir $\vb e_3$ 
	daugybos sąryšius.
	\begin{align}
		{\vb e_i}^2 &= \mathbb1			&\quad	{\gamma^i}^2 &= \mqty(\dmat{I,I})\\
		\vb e_i \vb e_j &= -\vb e_j \vb e_i	&\quad	(i\gamma^i)(i\gamma^j) &= - (i\gamma^j)(i\gamma^i)
	\end{align}.
	Matome, kad bazė atvaizduojama į tieškai neprikausomas matricas,
	tai turime ištikimą matricinį algebros $\Clp$ įvaizdį.
\end{proof}

Šį kartą lyginiai bazės elementai yra atviazduojami į blok-diagonamias matricas,
kurios yra dvi kopijos įvaiždžio iš teiginio \ref{prep:pauli_rep+} ir lyginio elemento bendra forma yra
\begin{equation}
		\mqty(\dmat{
			e^{\frac{i\theta}{2}(\vb n , \symbf\sigma)},
			e^{\frac{i\theta}{2}(\vb n , \symbf\sigma)}
		})
		=
		\exp(-\frac{\theta}{2}(
			 n_1\gamma^2\gamma^3
			+n_2\gamma^3\gamma^1
			+n_3\gamma^1\gamma^2
		))
		\,.
\end{equation}
Nelyginio elemeto bendra forma yra
\begin{equation}
		\mqty(\admat{
			 e^{\frac{i\theta}{2}(\vb n , \symbf\sigma)},
			-e^{\frac{i\theta}{2}(\vb n , \symbf\sigma)}
		})
		=
		-\gamma^1\gamma^2\gamma^3 \exp(-\frac{\theta}{2}(
			 n_1\gamma^2\gamma^3
			+n_2\gamma^3\gamma^1
			+n_3\gamma^1\gamma^2
		))
		\,.
\end{equation}










\section*{Išvados\addcontentsline{toc}{section}{Išvados}}
%\input{sec/conclusion}
\clearpage{}

\addcontentsline{toc}{section}{Literatūra}
\printbibliography

\clearpage{}
\begin{center}
	\textbf{ \large{
		THE DOUBLE GROUPS OF POINT GROUPS
		AND THEIR IRREDUCIBLE REPRESENTATIONS
	}} \large\par
\par\end{center}

\begin{center}
\textbf{Martynas Lašas}
\par\end{center}

\begin{center}
\textbf{Summary}
\par\end{center}

\input{sec/summary}


\end{document}
